\chapter{Analiza cerințelor}

\section{Introducere}

Înainte să încep să scriu cod, a trebuit să stabilesc cine folosește platforma și
ce trebuie să poată face fiecare. În capitolul ăsta trec prin tipurile de
utilizatori și prin cerințele care au ieșit din problema descrisă mai devreme.

\section{Actorii sistemului}

Platforma distinge patru tipuri de actori, ultimele trei reprezentate tehnic prin
roluri reale din \emph{ASP.NET Identity}: \textbf{Vizitator} (persoană
neautentificată, poate accesa pagina principală, se poate înregistra sau
autentifica), \textbf{Elev} (rol \texttt{Learn}, caută un profesor, trimite cereri
de colaborare, comunică prin chat cu profesorii cu care colaborează),
\textbf{Profesor} (rol \texttt{Teach}, poate fi găsit de elevi, acceptă sau
respinge cereri de colaborare, poate trimite fișiere elevilor cu care
colaborează) și \textbf{Administrator} (rol \texttt{Admin}, gestionează conturile,
blocările și materiile disponibile, moderează mesajele).

Fiecare cont primește, în plus, rolul implicit \texttt{User} la înregistrare.
Acesta nu influențează nicio funcționalitate vizibilă -- există exclusiv pentru
vizibilitate administrativă, ca un administrator să poată confirma din tabela de
roluri că un cont este un membru înregistrat obișnuit.

\section{Cerințe funcționale}

\begin{itemize}
    \item Înregistrare cu nume de utilizator, email și parolă; contul rămâne
    neconfirmat până la introducerea corectă a unui cod OTP trimis pe email.
    \item După confirmare, utilizatorul parcurge un flux de completare a profilului
    în trei pași: alegerea rolului, detaliile specifice rolului (vârstă, gen, nivel
    de studii/experiență) și selectarea a până la trei materii de interes.
    \item Căutarea altor utilizatori după nume, rol, nivel și materie.
    \item Trimiterea unei cereri de colaborare de la elev către profesor; acceptare,
    respingere sau anulare (cât timp e în așteptare) a cererii.
    \item Comunicare prin chat în timp real între utilizatorii care colaborează, cu
    posibilitatea profesorului de a atașa fișiere (documente, imagini, prezentări)
    conversației, și un indicator de mesaje necitite în bara de navigare.
    \item Editarea profilului propriu (biografie, avatar, materii, nivel); modificările
    sensibile (nume de utilizator, email) necesită reconfirmare prin cod OTP.
    \item Administrare: promovare/retrogradare a rolului de administrator, blocare
    temporară (1 zi, 7 zile, 1 lună) sau permanentă, ștergere de cont, gestionarea
    materiilor disponibile.
\end{itemize}

\section{Cerințe nefuncționale}

\begin{itemize}
    \item \textbf{Securitate}: parolele nu sunt stocate în clar (gestionate de
    ASP.NET Identity ca \emph{hash}); fișierele încărcate sunt validate după
    extensie și dimensiune; acțiunile administrative sunt restricționate exclusiv
    rolului \texttt{Admin}, verificat de framework, nu doar ascuns din interfață;
    toate cererile care modifică starea aplicației sunt însoțite de un token
    antiforgery validat pe server.
    \item \textbf{Localizare}: interfața este disponibilă integral în română și
    engleză, comutabilă fără reîncărcare manuală și fără autentificare.
    \item \textbf{Utilizabilitate}: interfața este responsivă și accesibilă atât
    pe desktop, cât și pe dispozitive mobile.
    \item \textbf{Extensibilitate}: arhitectura modulară (controllere separate
    pentru căutare, colaborări, cereri, chat, administrare) permite adăugarea de
    noi funcționalități fără a afecta modulele existente.
    \item \textbf{Fiabilitate}: acțiunile ireversibile (ștergere, blocare
    permanentă) sunt limitate prin verificări explicite -- un administrator nu
    își poate șterge sau bloca propriul cont, iar niciun administrator nu poate
    fi blocat sau șters de altul.
    \item \textbf{Mentenabilitate}: evoluția schemei bazei de date este gestionată
    exclusiv prin migrații Entity Framework Core, trasabile și reproductibile pe
    orice mediu nou.
\end{itemize}
